\SourceLeaf{201}
\(\mathrm{d}t\) und \(\mathrm{d}u\) lineare, Factoren zerlegt, muss entweder der eine oder der andere Factor \(=0\) werden, welches zwei verschiedene Integrationen geben wird. Die eine Integration wird der Gleichung
\[
\begin{aligned}
0={}&(aa+bb+cc)\mathrm{d}t\\
&+\bigl\{aa'+bb'+cc'+i\BareRoot{\bigl((aa+bb+cc)(a'a'+b'b'+c'c')}\\
&\qquad-(aa'+bb'+cc')^2\bigr)\bigr\}\mathrm{d}u
\end{aligned}
\]
\noindent entsprechen (wo \(i\) Kürze halber für \(\BareRoot{-1}\) geschrieben ist, indem man sich leicht überzeugt, dass der irrationale Theil des Ausdrucks imaginär werden muss); die andere einer ganz ähnlichen Gleichung, wenn nur \(i\) mit \(-i\) vertauscht wird. Ist also das Integral der erstern Gleichung dieses:
\[
p+iq=\text{Const.}
\]
\noindent wo \(p\) und \(q\) reelle Functionen von \(t\) und \(u\) bedeuten, so wird das andere Integral
\[
p-iq=\text{Const.}
\]
\noindent und die Natur der Sache wird es mit sich bringen, dass
\[
(\mathrm{d}p+i\,\mathrm{d}q)\mathord{.}(\mathrm{d}p-i\,\mathrm{d}q)\quad\text{oder}\quad\mathrm{d}p^2+\mathrm{d}q^2
\]
\noindent ein Factor von \(\omega\), oder
\[
\omega=n(\mathrm{d}p^2+\mathrm{d}q^2)
\]
\noindent werden muss, wo \(n\) eine endliche Function von \(t\) und \(u\) sein wird.

Wir wollen nun das Trinomium, in welches
\[
\mathrm{d}X^2+\mathrm{d}Y^2+\mathrm{d}Z^2
\]
\noindent übergeht, wenn für \(\mathrm{d}X,\mathrm{d}Y,\mathrm{d}Z\) ihre Werthe durch \(T,U,\mathrm{d}T,\mathrm{d}U\) substituirt werden, durch \(\Omega\) bezeichnen, und annehmen, dass auf ähnliche Weise, wie vorher, die beiden Integrale der Gleichung \(\Omega=0\) diese seien:
\[
\begin{aligned}
P+iQ&=\text{Const.}\\
P-iQ&=\text{Const.}
\end{aligned}
\]
\noindent und
\[
\Omega=N(\mathrm{d}P^2+\mathrm{d}Q^2)
\]
\noindent wo \(P,Q,N\) reelle Functionen von \(T\) und \(U\) bedeuten werden.
